\section{Systems Under Test}
\label{sec:systems}

Figure~\ref{fig:systems} shows the four engines side by side. All four serve an
OpenAI-compatible HTTP API, run one request at a time, and verify drafts exactly.

\begin{figure*}[t]
  \centering
  \begin{tikzpicture}[
      font=\footnotesize,
      box/.style={draw=blue!60!black, fill=blue!55!black, text=white, rounded corners=2pt,
                  minimum width=3.7cm, minimum height=0.75cm, align=center},
      weights/.style={draw=gray!70, fill=gray!10, rounded corners=2pt,
                      minimum width=3.7cm, minimum height=0.75cm, align=center},
      arr/.style={-{Stealth[length=2mm]}, gray!70!black, thick}]
    % columns
    \node[box] (l1) at (0,0) {llama-server\\ggml graph, CUDA graphs};
    \node[box] (l2) at (0,-1.25) {MTP head\\2 drafts};
    \node[weights] (l3) at (0,-2.5) {UD-Q4\_K\_XL GGUF\\17.6\,GB};
    \node[box] (o1) at (4.3,0) {OpenJet server\\one persistent kernel};
    \node[box] (o2) at (4.3,-1.25) {MTP head on device\\4 drafts, 32K vocab};
    \node[weights] (o3) at (4.3,-2.5) {same GGUF\\re-laid-out k-quants};
    \node[box] (n1) at (8.6,0) {ninfer-serve\\fused kernels, CUDA graphs};
    \node[box] (n2) at (8.6,-1.25) {MTP 3 drafts or\\DFlash2 7 drafts};
    \node[weights] (n3) at (8.6,-2.5) {.ninfer artifact\\4/5-bit groups, 16.7\,GiB};
    \node[box] (a1) at (12.9,0) {tqserve (azllm)\\fused kernels, prompt cache};
    \node[box] (a2) at (12.9,-1.25) {DFlash2 7 drafts\\0.79\,GB drafter};
    \node[weights] (a3) at (12.9,-2.5) {TQ\_L GGUF\\4.31 bits/weight, 14.7\,GB};
    \foreach \c in {l,o,n,a} { \draw[arr] (\c3) -- (\c2); \draw[arr] (\c2) -- (\c1); }
    \node[font=\small\bfseries] at (0,0.8) {Baseline};
    \node[font=\small\bfseries] at (4.3,0.8) {OpenJet};
    \node[font=\small\bfseries] at (8.6,0.8) {NInfer-4090};
    \node[font=\small\bfseries] at (12.9,0.8) {azllm};
  \end{tikzpicture}

  \caption{The four engines. Grey: the weights each one reads; blue: the drafter and the
           serving layer. OpenJet reads the baseline's own GGUF file; NInfer and azllm need their
           own weight files, both of which come with a DFlash2 drafter.}
  \label{fig:systems}
\end{figure*}

\paragraph{Baseline: llama.cpp with TurboQuant.}
The starting point is the TurboQuant fork of llama.cpp \citep{turboquant}, built from commit
\texttt{a3d5603d1} for \texttt{sm\_89}, serving Unsloth's \texttt{UD-Q4\_K\_XL} GGUF
 with flash attention, MTP speculation with two drafts, keys in
\texttt{q8\_0}, values in \texttt{turbo4} and a 128{,}000-token context. The vision projector
stays in host memory.

\paragraph{OpenJet Engine.}
OpenJet \citep{openjet} runs a whole speculative cycle (verification of $1+k$ tokens, lm-head,
acceptance, MTP catch-up and the next $k$ drafts) inside one cooperative kernel separated by
software grid barriers. It reads the same GGUF, keeps the k-quant encodings, and rolls the
DeltaNet state back by replaying accepted drafts from a small ring instead of snapshotting it.
The draft lm-head is restricted to the 32{,}768 most frequent tokens; verification uses the full
vocabulary. It is text only, and its KV cache is \texttt{f16} or \texttt{q8}.

\paragraph{NInfer-4090.}
NInfer \citep{ninfer} is a C++20/CUDA engine written for the Qwen3.6/3.8 family; we use the
\texttt{sm\_89} port \citep{ninfer4090}. It loads its own groupwise artifact
(4- and 5-bit weight groups, 8-bit output head) and offers MTP or DFlash2 drafting, a paged KV
cache in several formats (Table~\ref{tab:kvtypes}), prefix reuse at message boundaries, and an
optional vision tower. Its image targets CUDA~13.1; on the 580 driver (CUDA~13.0) it runs through
minor-version compatibility once the container's requirement check is disabled.

\paragraph{azllm (tqLLM).}
tqLLM \citep{tqllm} is a C++/CUDA engine for Qwen models stored in its own rotated
TurboQuant weight types, a GGUF extension that stock llama.cpp cannot load. azllm \citep{azllm}
is a fork that changes only the serving layer (streaming with tools, live timings, prompt-cache
bookmarks); we found it through a practitioner's report. We serve the
\texttt{TQ\_L} weights (4.31 bits per weight, 14.7\,GB, against about 5.2 for the other two
files) with a DFlash2 drafter quantized to TQ3 (0.79\,GB) and a 65{,}536-token
draft vocabulary. KV formats pair rotated 8- or 4-bit keys with 8-, 4-, 3- or 2-bit values; the
prompt cache keeps the previous request's KV so that a repeated prefix is not read again. The
vision projector must be F16 or F32. \texttt{TQ\_L} embeds Qwen's official chat template, which
fails on tool calls whose arguments arrive as a JSON string, so our launcher passes the fixed
template that the \texttt{TQ3\_4S-v2} file embeds.

\paragraph{Engines not retained.}
A Windows-oriented fork of NInfer-4090 has the same single-token decode rate
(about 52\,tok/s) and its KV modes are already merged into the port we use. vLLM and SGLang with
4-bit AWQ/GPTQ weights leave too little memory for a useful KV cache on 24\,GB, and NVFP4 needs
Blackwell. Lucebox runs DFlash2 inside a llama.cpp-based server but publishes
no RTX~4090 figure, and the one community report of DFlash2 in llama.cpp on this card we found
gives 73--75\,tok/s.
